\section{Privacy}

Anonymization (removing name, address, etc) is a typical
approach for private release. However, there are only
so many bits of anonimity before an individual is uniquely
identified. Something as simple as zip code, favorite
song, and age is enough to uniquely identify almost
anyone.

If stats are calculated over a small group of people,
it's easy to make inferences that identify those
people. One solution is $k$-anonimity, where each
individual's data is indistinguishable from at least
$k - 1$ other individuals' data.

Consider a curator who collects a sensitive bit $x_i$ from a
population of $n$ individuals.
To preserve privacy and allow the curator to accurately
infer aggregate statistics, each individual reports
\begin{equation}
    y =
    \begin{cases}
        x_i     & \text{with probability } \frac{1}{2} + \gamma \\
        1 - x_i & \text{with probability } \frac{1}{2} - \gamma
    \end{cases}
\end{equation}
for some fixed $\gamma \in [0, \frac{1}{2}]$.
This allows the curator to accurately estimate the mean
while still preserving privacy.

Differential Privacy is another approach, which requires
the chance that the released result will be $C$ is nearly
the same, whether or not you submit your information.

\begin{itemize}
    \item DP is closed under-post processing
    \item For a DP algorithm $M$, any function of the output of $M$
          is also differentially private
    \item DP algorithms can be composed. For other privacy mechanisms,
          releasing two algorithms that are individually
          private may completely break privacy.
    \item For DP, level of privacy degrades linearly with the number of times the private
          data is used
\end{itemize}

Noising the released data with the Laplace Distribution can also 
achieve differential privacy. 